%%^^A%%  fontspec-example.tex -- part of FONTSPEC <latex3.github.io/fontspec>

\documentclass{article}

\usepackage{fontspec}
\usepackage{xeCJK}
\setCJKmainfont{SourceHanSerifSC-Regular.otf}
\setCJKsansfont{SourceHanSansSC-Regular.otf}
\setCJKmonofont{SourceHanSansSC-Regular.otf}

\setmainfont{TeX Gyre Pagella}
\setsansfont{TeX Gyre Heros}[Scale=MatchLowercase]
\setmonofont{Inconsolata}[Scale=MatchLowercase]

\begin{document}
\pagestyle{empty}

\section*{\textsf{fontspec} 宏包基础}

\textsf{fontspec} 宏包为使用 Xe\TeX{} 或 Lua\TeX{} 排版的
\LaTeX{} 文档实现自动字体选择。基本命令为

{\centering \verb|\fontspec{font display name}[font features]|。\par}

默认字体、无衬线字体与打字机字体可分别用命令
\verb|\setmainfont|、\verb|\setsansfont| 与 \verb|\setmonofont| 设置，
如导言区所示。它们的语法与 \verb|\fontspec| 命令相同。
所有预期的字体字形均可用：

\begin{center}
  {\itshape Italics and \scshape small caps\dots}\\
  {\sffamily\bfseries Bold sans serif and \itshape bold italic sans serif\dots}
\end{center}

数学模式中的文字字体也会随之改变（例如注意 `$\cos(n\pi)=\pm 1$'
中的余弦函数），但前提是罗马字体与无衬线字体已在导言区设置；
在文档中途调用 \verb|\setmainfont| 不会影响这些数学模式字体。
数学符号本身不受影响。

注意导言区加载默认字体时所用的字体特性。罗马字体与无衬线字体会自动启用
\verb|Ligatures=TeX|，从而支持 \verb|``---''| 这类常规 \TeX{} 连字，
得到“---”。\verb|Scale=MatchLowercase| 会自动把各字体缩放到相同的 x 高。

更多信息请参阅 \textsf{fontspec} 完整文档。

\end{document}
